\usepackage{newunicodechar}
\usepackage{amsmath,amssymb}

% ── Emoji ────────────────────────────────────────────────────────────────────
% 📋 U+1F4CB  clipboard icon (decorative bullet in headings)
\newunicodechar{📋}{\,\textbullet\,}

% ── Math symbols written outside math mode ───────────────────────────────────
\newunicodechar{∅}{$\emptyset$}
\newunicodechar{≳}{$\gtrsim$}
\newunicodechar{⍵}{$\omega$}

% ᐧ U+1427  Canadian syllabics middle dot → multiplication dot
\newunicodechar{ᐧ}{$\cdot$}

% Angle brackets – two different Unicode pairs appear in the sources
\newunicodechar{〈}{$\langle$}   % U+3008  CJK left  angle bracket
\newunicodechar{〉}{$\rangle$}   % U+3009  CJK right angle bracket
\newunicodechar{⟨}{$\langle$}   % U+2329  left-pointing  angle bracket
\newunicodechar{⟩}{$\rangle$}   % U+232A  right-pointing angle bracket

% Ⲫ U+2CAA  Coptic Fi → used as Φ
\newunicodechar{Ⲫ}{$\Phi$}

% ── Mathematical italic letters used inline ───────────────────────────────────
\newunicodechar{𝑁}{$N$}   % U+1D441
\newunicodechar{𝑃}{$P$}   % U+1D443
\newunicodechar{𝑇}{$T$}   % U+1D447
\newunicodechar{𝑀}{$M$}   % U+1D440
\newunicodechar{𝑚}{$m$}   % U+1D45A
\newunicodechar{𝑛}{$n$}   % U+1D45B

% ── Mathematical bold-italic Greek ───────────────────────────────────────────
\newunicodechar{𝜶}{$\alpha$}   % U+1D736
\newunicodechar{𝜷}{$\beta$}    % U+1D737